% ============================================================
%  articles/a5-agreement.tex
%  Reaching the Impossible Agreement — Taha Hbirri
% ============================================================

\smjarticle{art:agreement}{%
  title       = {Reaching the Impossible Agreement},
  % cover-title = {},
  % toc-title   = {},
  author      = {Taha Hbirri},
  year        = {CS'27},
  affil       = {UgS V},
  category    = {article},
  abstract    = {Reaching agreement in a distributed system is a fundamental problem in
    computer science. We study the consensus problem, in which each process proposes a
    value and all correct processes must decide on a common value satisfying agreement,
    validity, and termination. The central result of this paper is the
    Fischer--Lynch--Paterson (FLP) impossibility theorem, which states that no
    deterministic algorithm can solve consensus in an asynchronous message-passing
    system tolerating even a single crash failure. We present the full proof. We then
    study Paxos which, under partial synchrony, can achieve consensus through a
    two-phase proposal protocol followed by a learning phase; its safety guarantee
    reduces to a single combinatorial invariant. Notably, this safety guarantee holds
    even without partial synchrony, which is needed only to guarantee that the protocol
    eventually terminates. We prove this safety property.},
}

\import{articles/papers/agreement/}{main.tex}